\subsection{Future Directions: Computational Scaling to Numerical Models}
\label{sec:future:models}

The metric-consistent spectral framework developed herein for localized observations is structurally extensible to high-resolution numerical simulations, such as Large Eddy Simulations (LES) or the Weather Research and Forecasting (WRF) model. Computing $D_{\mathrm{eff}}$ and $F_W$ across a full 3D spatial grid would unlock new capabilities in boundary layer analysis, including the spatial tracking of moving intermittent regimes and the diagnostic monitoring of Sub-Grid Scale (SGS) energy transfers.

However, scaling this framework to model outputs introduces a fundamental grid-topology mismatch. Pseudospectral expansions are strictly optimized for specific coordinate compactifications (e.g., the hyperbolic coordinate stretching applied to the CASES-99 tower domain), whereas numerical models typically utilize logarithmic, grid-nested, or terrain-following $\eta$-coordinate systems. Directly projecting model grids onto the tower's Chebyshev basis without adjustment would introduce severe numerical artifacts, such as Runge's phenomenon or spectral ringing at the domain boundaries.

Consequently, adapting this framework for model validation requires one of two strategic pathways: (1) establishing a 'virtual tower' array wherein dense model columns are vertically interpolated onto the established instrument coordinate nodes prior to matrix decomposition, or (2) re-deriving an independent \emph{UnifiedManifoldWorkspace} where the quadrature weights and mass matrix $\mathbf{M}$ are explicitly recomputed to match the native discrete grid-spacing of the numerical solver. The execution of these model-grid transformation pathways is reserved for a follow-up validation study.

\textbf{Applicability to Numerical Model Validation:}
The mathematical architecture of the $D_{\mathrm{eff}}$ pipeline is designed with explicit forward extensibility for numerical model evaluation. Rather than validating simulations purely through bulk, time-averaged statistics (e.g., mean potential temperature or surface fluxes), this framework introduces a structural complexity metric based on scale-by-scale energy decomposition. This capability provides an objective constraint for evaluating how models represent intermittent turbulence and wave-turbulence transitions. The underlying Julia package structure isolates the core Riemannian projection mathematics from the input grid data, ensuring compatibility with virtual model profiles once grid-mapping transformations are applied.

